%! AP Statistics — Unit 1, Lesson 1.4 — Group Activity ANSWER KEY
\documentclass[10pt]{article}
\usepackage{apstats-article}
\usepackage{apstats-key}

%=============================================================================
\begin{document}

\pageheader{Unit 1, Lesson 1.4}{Group Activity --- Answer Key}

\vspace{0.1in}

\begin{tcolorbox}[colback=sky, colframe=navy, boxrule=0.8pt, arc=2mm,
    left=3mm, right=3mm, top=2mm, bottom=2mm]
\small
\textbf{Part 1} (group, 8 min): Complete the table and construct a relative
frequency bar chart. Then write an interpretation.\\[4pt]
\textbf{Part 2} (group, 5 min): Identify and describe a misleading version.
Optional: compute sector angles for a pie chart.\\[4pt]
\textbf{Part 3} (individual, 3 min): Reflection question.
\end{tcolorbox}

\vspace{0.14in}

% ── PART 1 ───────────────────────────────────────────────────────────────────
{\large\bfseries\color{navy} Part 1 \quad Construct the Bar Chart}

\vspace{0.1in}

\noindent\small A survey of \textbf{50 Statistics students} asks: ``Which streaming
service do you use most?'' Results are below.

\vspace{0.1in}

\renewcommand{\arraystretch}{2.0}
\begin{tabularx}{\linewidth}{@{} >{\bfseries}p{0.22\linewidth}
                                    C{0.18\linewidth}
                                    C{0.22\linewidth}
                                    C{0.22\linewidth} @{}}
\rowcolor{sky}
\textbf{Service} & \textbf{Frequency} & \textbf{Rel.\ Frequency} & \textbf{Angle (${}^\circ$)} \\
\midrule
Netflix  & 21 & \ans{0.42 (42\%)} & \ans{151.2} \\
Disney+  & 10 & \ans{0.20 (20\%)} & \ans{72}    \\
Hulu     & 9  & \ans{0.18 (18\%)} & \ans{64.8}  \\
HBO Max  & 6  & \ans{0.12 (12\%)} & \ans{43.2}  \\
Other    & 4  & \ans{0.08 (8\%)}  & \ans{28.8}  \\
\midrule
\rowcolor{sky}\textbf{Total} & \ans{50} & \ans{1.00} & \ans{360} \\
\end{tabularx}

\vspace{0.14in}

\noindent\textbf{Construct a relative frequency bar chart below.}\\[2pt]
\small Include: (1) a title, (2) a labeled x-axis with categories, (3) a labeled
y-axis starting at 0, (4) gaps between bars.

\vspace{8pt}
\begin{tcolorbox}[colback=white, colframe=navy, boxrule=0.8pt, arc=2mm,
    left=3mm, right=3mm, top=2mm, bottom=2mm, height=2.8in]
\centering
{\color{keyred}\bfseries\small Preferred Streaming Service, $n=50$}\\[2pt]
\begin{tikzpicture}[scale=1.0]
  % y-axis 0 to 0.45, 1 unit = 10cm  →  0.45 = 4.5cm
  \draw[thick] (0,0) -- (0,4.7);
  \foreach \p/\lbl in {0/0, 1/0.10, 2/0.20, 3/0.30, 4/0.40} {
    \draw (-0.1,\p) -- (0.1,\p);
    \node[left, font=\scriptsize] at (-0.12,\p) {\lbl};
  }
  \node[rotate=90, font=\scriptsize, color=keyred] at (-1.05,2.3) {Relative frequency};
  \draw[thick] (0,0) -- (8.2,0);
  % bars: value*10 cm tall, width 1.0, gap 0.6
  \foreach \i/\v/\lbl in {0/4.2/Netflix, 1/2.0/Disney+, 2/1.8/Hulu,
                          3/1.2/{HBO Max}, 4/0.8/Other} {
    \fill[keyred!75] ({0.6+\i*1.5},0) rectangle ({1.6+\i*1.5},\v);
    \draw[keyred]    ({0.6+\i*1.5},0) rectangle ({1.6+\i*1.5},\v);
    \node[below, font=\scriptsize, color=keyred] at ({1.1+\i*1.5},-0.08) {\lbl};
  }
  \node[below, font=\scriptsize\itshape, color=keyred] at (4.1,-0.55) {Streaming service};
\end{tikzpicture}
\end{tcolorbox}

\vspace{0.12in}

\begin{enumerate}[leftmargin=1.5em, itemsep=8pt]
  \item Does your relative frequency column sum to 1? \ans{Yes}
        If not, find and fix the rounding error.

  \item Write a complete sentence identifying the most popular and least popular
        streaming service, including relative frequencies and context.\\[5pt]
        \ansline{Of the 50 Statistics students surveyed, Netflix was the most}\\[1pt]
        \ansline{popular service (42\%) and ``Other'' the least popular (8\%).}

  \item Write a second sentence describing the overall pattern of the distribution
        (e.g., is one service much more popular than the others?).\\[5pt]
        \ansline{The distribution is dominated by Netflix --- no other service reaches}\\[1pt]
        \ansline{half its share; Disney+ (20\%) and Hulu (18\%) are close behind.}
\end{enumerate}

\vspace{0.14in}

% ── PART 2 ───────────────────────────────────────────────────────────────────
{\large\bfseries\color{navy} Part 2 \quad Misleading Graphs \& Pie Chart (Optional)}

\vspace{0.1in}

\begin{enumerate}[leftmargin=1.5em, itemsep=8pt]

  \item Suppose someone redraws your bar chart with the y-axis starting at 30\%
        instead of 0\%. Describe how this changes what a reader sees. Which service
        would appear most dramatically more popular?\\[5pt]
        \ansline{Only Netflix (42\%) would clear the 30\% cutoff --- every other bar}\\[1pt]
        \ansline{vanishes below the axis. Netflix would look near-exclusive, when in}\\[1pt]
        \ansline{fact the other services account for 58\% of responses combined.}

  \item A student presents a 3-D pie chart of this data with the Netflix slice
        tilted toward the viewer. Explain why this display might be misleading
        even though the data are correct.\\[5pt]
        \ansline{Perspective enlarges the near slice, so readers overestimate Netflix's}\\[1pt]
        \ansline{42\% share. The data are right but the presentation inflates it.}

  \item \textbf{Optional --- Pie Chart:} Using the angles you computed in the table,
        describe the largest and smallest sectors. Is a bar chart or pie chart more
        effective for this dataset? Why?\\[5pt]
        \ansline{Largest: Netflix, $151.2^\circ$. Smallest: Other, $28.8^\circ$.}\\[1pt]
        \ansline{The bar chart wins --- comparing heights beats estimating angles.}

\end{enumerate}

\vspace{0.14in}

% ── PART 3 ───────────────────────────────────────────────────────────────────
{\large\bfseries\color{navy} Part 3 \quad Individual Reflection}

\vspace{0.1in}

\begin{tcolorbox}[colback=goldbg, colframe=goldacc, boxrule=0.8pt, arc=2mm,
    left=3mm, right=3mm, top=2mm, bottom=2mm]
\small\textbf{In your own words:} What is the single most important thing to check
when you look at a bar chart made by someone else? Why does it matter?
Give a specific example from today's activity.\\[6pt]
\ansline{Check that the y-axis starts at 0. If it starts higher, differences between}\\[1pt]
\ansline{bars are visually exaggerated --- starting at 30\% here would make Netflix}\\[1pt]
\ansline{look dominant while hiding that Disney+ and Hulu are substantial.}
\end{tcolorbox}

\end{document}
